\chapter{Computation Time Study}
\label{ch:performance}

RVE data generation is the dominant offline cost because each sample requires
mixed force equilibrium, rotational equilibrium, homogenization, and tangent
condensation. The implementation therefore supports atomic checkpoints and
parallel sample generation.

\begin{table}[H]
\centering
\caption{Available performance evidence and quantities still unmeasured.}
\label{tab:performance}
\begin{tabularx}{\textwidth}{@{}Xr@{}}
\toprule
Quantity & Value or status\\
\midrule
Mean nonlinear RVE solve time & Not recorded\\
Mean tangent condensation time & Not recorded\\
One-worker elapsed time & \(1709.97\,\mathrm{s}\)\\
Twelve-worker elapsed time & \(513.22\,\mathrm{s}\)\\
Parallel speedup and efficiency & \(3.332\times\), \(27.8\%\)\\
Recorded maximum RSS, 1/12 workers & 301460/297972 kB\\
Surrogate stress/tangent evaluation time & Not recorded\\
Direct-RVE/surrogate speedup & Not measured\\
\bottomrule
\end{tabularx}
\end{table}

The two timing files record the same benchmark campaign but do not preserve
enough command metadata to derive a defensible per-sample throughput. Their
elapsed-time ratio gives the speedup above. The low parallel efficiency can
include process startup, serial assembly, unequal sample difficulty, and
file-system contention; the records do not isolate these causes. Likewise,
the maximum-RSS accounting scope is not documented well enough to call the
12-worker figure ``per worker.'' A separate attempted parallel run terminated
by signal~2 after \(452.72\,\mathrm{s}\) and is excluded from speedup claims.
No direct-RVE/surrogate timing was captured, so those quantities are reported
as unavailable rather than estimated.
